\documentclass[12pt]{article}
\usepackage{amsmath, amssymb, amsthm, bm}
\usepackage[margin=1in]{geometry}
\usepackage[skip=0.6em]{parskip}

\newtheorem{proposition}{Proposition}
\newtheorem{lemma}{Lemma}
\theoremstyle{remark}
\newtheorem{remark}{Remark}

\begin{document}

\begin{center}
{\Large The three-dimensional core--redundant case}\\[3pt]
{\small Ivan A. Canay :: August 28, 2026}
\end{center}

\medskip

Consider
\[
X\sim N_3(\mu,I_3)~,
\]
where coordinates 1 and 2 are core moments and coordinate 3 is redundant.  The
maintained mean space is
\[
\Lambda
:=
\left\{\mu\in\mathbf R^3:
\mu_3\le\max\{\mu_1,\mu_2\}
\right\}~.
\]
The null and alternative are
\[
\Lambda_0:=\Lambda\cap(-\infty,0]^3,
\qquad
\Lambda_1
:=
\left\{\mu\in\Lambda:
\max\{\mu_1,\mu_2\}>0
\right\}~.
\]
Fix $0<\alpha<1$.
A level-$\alpha$ test is called $\alpha$-admissible if there is no other
level-$\alpha$ test whose power is weakly larger at every
$\mu\in\Lambda_1$ and strictly larger at some $\mu\in\Lambda_1$.

The test of primary interest is the equal-cutoff rectangle, or max test.  Its
acceptance region is
\begin{equation}\label{eq:rectangle}
A_\kappa
:=
\bigcap_{j=1}^3\{x_j\le\kappa\}~,
\end{equation}
and its rejection rule is
\[
\phi_{\max}(x)
=
I\left\{\max_{1\le j\le3}x_j>\kappa\right\}~.
\]
Thus the statistic is the maximum of all three coordinates and a single
critical value is applied to that maximum.  Under independence, exact size at
the origin requires
\[
\Phi(\kappa)^3=1-\alpha,
\qquad
\kappa=\Phi^{-1}\{(1-\alpha)^{1/3}\}~.
\]
Coordinatewise lowerness then makes the test level $\alpha$ over
$\Lambda_0$, with size attained at $\mu=0$.

\section{Failure of the feasible-normal sufficient condition}

For the present identity-covariance experiment, the feasible-normal set in
Proposition~1 of the extension plan becomes
\begin{equation}\label{eq:feasible-set}
\mathcal F(I_3,\Lambda)
=
\left\{a\in\mathbf R^3:
a\ge0,\quad a_3\le\max\{a_1,a_2\}
\right\}~.
\end{equation}
Indeed, the general condition is $a\ge0$ and $I_3a=a\in\Lambda$.

The support function of the rectangle in~\eqref{eq:rectangle} is
\[
h_{A_\kappa}(a)
=
\begin{cases}
\kappa(a_1+a_2+a_3),&a\ge0,\\
+\infty,&\text{otherwise}.
\end{cases}
\]
Consequently, every nonzero direction $a\ge0$ is a finite-support direction
of $A_\kappa$.  In particular, the normal to the redundant-coordinate face is
\[
e_3=(0,0,1),
\qquad
h_{A_\kappa}(e_3)=\kappa<\infty~.
\]
But
\[
(e_3)_3=1
>
\max\{(e_3)_1,(e_3)_2\}=0~,
\]
so $e_3\notin\mathcal F(I_3,\Lambda)$.  The equal-cutoff rectangle therefore
fails Proposition~1 of the extension plan.

The same conclusion holds for every finite three-dimensional rectangle, not
only for the equal-cutoff rectangle, because each such rectangle has the
finite normal $e_3$.  This calculation does not establish that the max test is
inadmissible.  The extension-plan result is sufficient, not necessary, and
its failure leaves the admissibility of the rectangle unresolved.

\section{Limits of the two-dimensional gated repair}

In the two-dimensional problem with one fixed core moment, a rectangular face
for the redundant moment can be replaced by a calibrated gate involving the
sum of the core and redundant statistics.  The power comparison works because
every alternative satisfies both
\[
\mu_c>0
\qquad\text{and}\qquad
\mu_3\le\mu_c
\]
for the same fixed core coordinate $c$.

There is no such fixed coordinate in the present model.  The alternative is
the union of the two branches
\[
\Lambda_1^{(1)}
:=
\{\mu:\mu_1>0,\ \mu_3\le\mu_1\},
\qquad
\Lambda_1^{(2)}
:=
\{\mu:\mu_2>0,\ \mu_3\le\mu_2\}~,
\]
where the unpaired core mean is unrestricted on each branch.  On the first
branch, the likelihood-ratio exponent can be written as
\begin{equation}\label{eq:branch-one-lr}
\mu'x
=
\mu_1(x_1+x_3)
+(\mu_1-\mu_3)(-x_3)
+\mu_2x_2~.
\end{equation}
The first two coefficients in~\eqref{eq:branch-one-lr} are nonnegative in the
directions required by the bivariate transfer argument.  This supports a gate
pairing coordinates 1 and 3, while coordinate 2 is left unchanged.  The
analogous decomposition on the second branch supports a gate pairing
coordinates 2 and 3.  A gate based on one fixed core coordinate need not
improve power on the other branch: that selected core mean may be negative and
may lie below $\mu_3$.

A natural attempt to avoid selecting a core coordinate in advance is the
symmetric gated acceptance region
\begin{equation}\label{eq:symmetric-gate}
A_b^{\mathrm g}
:=
\{x_1\le\kappa,\ x_2\le\kappa,\
  x_1+x_3\le b,\ x_2+x_3\le b\}~.
\end{equation}
Equivalently, the redundant face is replaced by
\[
x_3+\max\{x_1,x_2\}\le b~.
\]
Choose $b$ so that
\[
P_0(A_b^{\mathrm g})=P_0(A_\kappa)=1-\alpha~.
\]
The finite facet normals of $A_b^{\mathrm g}$ are
\[
e_1,\quad e_2,\quad (1,0,1),\quad (0,1,1)~,
\]
and each belongs to~\eqref{eq:feasible-set}.  This alone is not enough for the
extension-plan condition because the feasible-normal set is not convex and
that condition concerns every finite-support direction.  In particular, the
sum of the two gated-face
normals is $(1,1,2)$.  It has
$h_{A_b^{\mathrm g}}(1,1,2)=2b<\infty$, but is not feasible because
$2>\max\{1,1\}$.  Thus the symmetric gated test, like the original rectangle,
falls outside the sufficient condition.  More importantly here, it does not
dominate the original rectangle.

In fact, the symmetric gate does not dominate.  Its acceptance probability at
a general mean is
\begin{equation}\label{eq:gated-acceptance-probability}
P_\mu(A_b^{\mathrm g})
=
\int_{-\infty}^{\infty}
\varphi(z-\mu_3)
\Phi\!\left(\min\{\kappa,b-z\}-\mu_1\right)
\Phi\!\left(\min\{\kappa,b-z\}-\mu_2\right)
\,dz~,
\end{equation}
whereas
\begin{equation}\label{eq:rectangular-acceptance-probability}
P_\mu(A_\kappa)
=
\prod_{j=1}^3\Phi(\kappa-\mu_j)~.
\end{equation}
For example, take $\alpha=0.05$.  Null calibration gives
\[
\kappa\simeq2.121201,
\qquad
b\simeq3.117267~.
\]
At
\[
\mu=(0.1,-5,0.1)\in\Lambda_1~,
\]
equations~\eqref{eq:gated-acceptance-probability} and
\eqref{eq:rectangular-acceptance-probability} give
\[
P_\mu(A_\kappa)\simeq0.957209,
\qquad
P_\mu(A_b^{\mathrm g})\simeq0.965438~.
\]
The corresponding rejection probabilities are approximately $0.042791$ and
$0.034562$.  The symmetric gated test therefore has lower power at this
alternative.

The example exposes the adaptation cost created by the maximum restriction.
When $\mu_2$ is very negative, coordinate 1 is effectively the only relevant
core coordinate.  Nevertheless, the symmetric gate is calibrated under the
null to allow either core coordinate to activate the gate.  The resulting
critical value is too large in this nearly bivariate subproblem.

There is also a geometric obstruction to a simple symmetrization.  A
sample-space move $d$ that increases every alternative likelihood ratio would
have to satisfy
\[
\mu'd\ge0
\qquad\text{for every }\mu\in\Lambda_1~.
\]
The unrestricted value of $\mu_2$ on $\Lambda_1^{(1)}$ forces $d_2=0$, and
the unrestricted value of $\mu_1$ on $\Lambda_1^{(2)}$ forces $d_1=0$.
With these equalities, taking $\mu_3=\mu_1>0$ on the first branch gives
$d_3\ge0$, while allowing $\mu_3\to-\infty$ gives $d_3\le0$.  Hence the
common likelihood-ratio transfer cone is
\[
\{0\}~.
\]
There is no nonzero pointwise move that improves every alternative in both
branches.  Any dominating test, if one exists, must therefore use a genuinely
distributional argument that adapts to the information in the data about
which core coordinate is active.  This observation rules out the direct
transfer argument, but it does not determine admissibility.  The next two
sections give sharper sufficient conditions that establish admissibility of
both the symmetric gate and the rectangle.

\section{A sharper sufficient condition based on defining supports}

The sufficient condition in Section~1 requires every direction in which the
acceptance region has finite support to be feasible.  Fabel's observation is
that this is stronger than needed in the present two-core model.  It is enough
that the feasible supporting halfspaces already determine the acceptance
region.

Call an acceptance region $A$ \emph{nice} if it is closed, convex,
coordinatewise lower, and satisfies $P_0(A)=1-\alpha$.  Let
\[
\mathcal F
:=
\{a\in\mathbf R^3:a\ge0,\ a_3\le\max\{a_1,a_2\}\}
\]
and define
\[
F_A:=\{a\in\mathcal F\setminus\{0\}:h_A(a)<\infty\}~.
\]
We say that $A$ is \emph{feasibly supported} if
\begin{equation}\label{eq:feasibly-supported}
A
=
\bigcap_{a\in F_A}\{x:a'x\le h_A(a)\}~.
\end{equation}
For a polyhedral region, this means that its essential facet normals can all
be chosen feasible.

\begin{proposition}[Feasible support suffices]
\label{prop:feasible-support}
Every feasibly supported nice test is $\alpha$-admissible against all
level-$\alpha$ tests, including randomized competitors.
\end{proposition}

\begin{proof}
Let $\phi_A:=I\{X\notin A\}$, suppose that a randomized level-$\alpha$
test $\psi$ weakly dominates $\phi_A$, and put $h:=\psi-\phi_A$.  Since
$0\in\Lambda_0$, level control gives $E_0\psi\le\alpha=E_0\phi_A$.
Dominance along $t e_1\in\Lambda_1$ and continuity as $t\downarrow0$ give
the reverse inequality, so
\begin{equation}\label{eq:nice-zero-mass}
E_0h=0~.
\end{equation}

Fix $a\in F_A$.  Since $a$ is nonzero and feasible, $a\in\Lambda_1$.
Dominance along the ray $ta$, together with the Gaussian likelihood ratio,
gives
\begin{equation}\label{eq:feasible-tilt}
E_0[e^{t a'X}h(X)]\ge0
\qquad\text{for every }t>0~.
\end{equation}
Write $h=h^+-h^-$.  The sign restrictions imposed by $\phi_A$ give
$\{h^+>0\}\subseteq A$.  If, for some $\varepsilon>0$,
\[
m_{a,\varepsilon}
:=
E_0[h^-1\{a'X\ge h_A(a)+\varepsilon\}]>0~,
\]
then
\[
E_0[e^{t a'X}h(X)]
\le
e^{t h_A(a)}E_0h^+
-
e^{t\{h_A(a)+\varepsilon\}}m_{a,\varepsilon}<0
\]
for all sufficiently large $t$, contradicting~\eqref{eq:feasible-tilt}.
Taking $\varepsilon=1/n$ shows that $h^-=0$ almost surely on
$\{a'X>h_A(a)\}$.

By~\eqref{eq:feasibly-supported}, the open set $A^c$ is the union of these
open halfspaces over $a\in F_A$.  The Lindel\"of property supplies a
countable subcover, so the preceding conclusion gives $h^-=0$ almost surely
on $A^c$.  But $\{h^->0\}\subseteq A^c$, hence $h^-=0$ almost surely.
Equation~\eqref{eq:nice-zero-mass} then forces $h^+=0$ almost surely.
Mutual absolute continuity of Gaussian shifts rules out a strict power
improvement at any alternative.
\end{proof}

The logical relation with the earlier sufficient condition is
\[
\left\{a\ge0:h_A(a)<\infty\right\}\subseteq\mathcal F
\quad\Longrightarrow\quad
\text{$A$ is feasibly supported}
\quad\Longrightarrow\quad
\text{$\phi_A$ is admissible}~.
\]
The first implication is not reversible.  The old condition controls every
finite-support direction, whereas Proposition~\ref{prop:feasible-support}
uses only the supporting halfspaces needed to reconstruct $A$.

\begin{remark}[The symmetric gate is admissible]
The four facet normals of $A_b^{\mathrm g}$ in~\eqref{eq:symmetric-gate} are
$e_1$, $e_2$, $(1,0,1)$, and $(0,1,1)$, all feasible.  Hence the symmetric
gate is admissible by Proposition~\ref{prop:feasible-support}.  Its bad
finite-support direction $(1,1,2)$ is a sum of facet normals and is not needed
to define the region.  Thus the symmetric gate and the rectangle are both
admissible, while the numerical calculation in Section~2 shows that the
symmetric gate does not dominate the rectangle.
\end{remark}

\section{Admissibility of the rectangular max test}

The failure of the feasible-normal condition does not imply inadmissibility in
the present model.  In fact, the equal-cutoff max test is admissible.  The
argument does not use equality of the cutoffs and therefore applies to every
finite exact-size rectangle.

For finite $c_1,c_2,c_3$, define
\begin{equation}\label{eq:general-rectangle}
A_c
:=
\{x_1\le c_1,\ x_2\le c_2,\ x_3\le c_3\},
\qquad
\phi_c:=I\{X\notin A_c\}~.
\end{equation}

\begin{proposition}[Admissibility of finite rectangles]
\label{prop:rectangle-admissibility}
Suppose that $P_0(A_c)=1-\alpha$.  Then $\phi_c$ is
$\alpha$-admissible against all level-$\alpha$ tests, including randomized
competitors.  In particular, the equal-cutoff max test in
\eqref{eq:rectangle} is $\alpha$-admissible.  More strongly, any
level-$\alpha$ test that weakly dominates $\phi_c$ coincides with it almost
surely under every $P_\mu$.
\end{proposition}

\begin{proof}
Suppose that a randomized level-$\alpha$ test $\psi$ weakly dominates
$\phi_c$ on $\Lambda_1$.  Define
\[
h(x):=\psi(x)-\phi_c(x)~.
\]
Because $\mu=0$ belongs to the null, level control gives
$E_0\psi\le\alpha=E_0\phi_c$.  On the other hand, $t e_1\in\Lambda_1$ for
every $t>0$.  Dominance along this ray and continuity as $t\downarrow0$ give
$E_0\psi\ge E_0\phi_c$.  Consequently,
\begin{equation}\label{eq:zero-mass}
E_0h=0~.
\end{equation}

We first show the first-moment identities $E_0[X_jh(X)]=0$, $j=1,2,3$.  Let
\[
q:=E_0[Xh(X)]\in\mathbf R^3~.
\]
The set $\Lambda_1$ is a cone, so $td\in\Lambda_1$ whenever
$d\in\Lambda_1$ and $t>0$.  Dominance and the Gaussian likelihood ratio
therefore imply
\[
0\le E_{td}h(X)
=
\exp\{-t^2\|d\|^2/2\}
E_0\!\left[h(X)\exp\{t d'X\}\right]~.
\]
The expression equals zero at $t=0$ by~\eqref{eq:zero-mass}.
Because $|h|\le1$ and the Gaussian moment-generating function is finite
everywhere, differentiation under the expectation is justified in a
neighborhood of zero.  Differentiating from the right at zero therefore gives
\begin{equation}\label{eq:directional-moment}
d'q\ge0
\qquad\text{for every }d\in\Lambda_1~.
\end{equation}
For every $s\in\mathbf R$, both $(1,s,0)$ and $(s,1,0)$ belong to
$\Lambda_1$.  Hence~\eqref{eq:directional-moment} gives
\[
q_1+s q_2\ge0,
\qquad
s q_1+q_2\ge0
\qquad\text{for every }s\in\mathbf R~,
\]
which forces $q_1=q_2=0$.  Moreover, $(1,0,r)\in\Lambda_1$ for every
$r\le1$.  Applying~\eqref{eq:directional-moment} with $r=1$ and $r=-1$
then gives $q_3=0$.  We have shown that
\begin{equation}\label{eq:zero-first-moments}
E_0[X_jh(X)]=0,
\qquad j=1,2,3~.
\end{equation}

It remains to use the sign restrictions imposed by the rectangular test.
Write $h=h^+-h^-$ for the positive and negative parts of $h$.  Since
$\phi_c=0$ on $A_c$ and $\phi_c=1$ outside $A_c$,
\begin{equation}\label{eq:sign-support}
\{h^+>0\}\subseteq A_c,
\qquad
\{h^->0\}\subseteq A_c^c~,
\end{equation}
pointwise.  Along the alternative ray $t e_j$, for $j=1,2$, dominance
gives
\begin{equation}\label{eq:laplace-nonnegative}
E_0\!\left[e^{tX_j}h(X)\right]\ge0
\qquad\text{for every }t>0~.
\end{equation}
We claim that
\begin{equation}\label{eq:no-core-loss}
h^-=0\ \text{a.s. on }\{X_j>c_j\},
\qquad j=1,2~.
\end{equation}
To see this, suppose for some $\varepsilon>0$ that
\[
m_{j,\varepsilon}
:=
E_0\!\left[h^-1\{X_j\ge c_j+\varepsilon\}\right]>0~.
\]
Let $M^+:=E_0h^+$.  The first inclusion in~\eqref{eq:sign-support} and
the preceding display imply
\[
E_0\!\left[e^{tX_j}h(X)\right]
\le
e^{tc_j}M^+
-e^{t(c_j+\varepsilon)}m_{j,\varepsilon}~,
\]
which is negative for all sufficiently large $t$.  This contradicts
\eqref{eq:laplace-nonnegative}.  Applying the argument with
$\varepsilon=1/n$, for $n\ge1$, proves~\eqref{eq:no-core-loss}.

By~\eqref{eq:sign-support} and~\eqref{eq:no-core-loss}, any negative part of
$h$ must be supported on
\[
\{X_1\le c_1,\ X_2\le c_2,\ X_3>c_3\}~,
\]
whereas the positive part is supported on $\{X_3\le c_3\}$.  Equation
\eqref{eq:zero-mass} gives $E_0h^+=E_0h^-=:M$.  If $M>0$, these support
restrictions imply
\[
E_0[X_3h^+]\le c_3M,
\qquad
E_0[X_3h^-]>c_3M~.
\]
The second inequality is strict because $X_3>c_3$ on $\{h^->0\}$ and
$E_0h^-=M>0$.  Consequently, $E_0[X_3h]<0$.
This contradicts~\eqref{eq:zero-first-moments}.  Thus $M=0$ and $h=0$
$P_0$-almost surely.  All Gaussian shifts are mutually absolutely continuous,
so $\psi$ and $\phi_c$ coincide almost surely under every $P_\mu$.  A strict
power improvement is therefore impossible.
\end{proof}

\begin{remark}[Why the same proof fails with one core]
In the two-dimensional model with one core and one redundant coordinate, the
alternative directions can be normalized as $(1,s)$ with $s\le1$.  If
$q=E_0[Xh(X)]$, the local dominance condition becomes
\[
q_1+s q_2\ge0
\qquad\text{for every }s\le1~,
\]
which is equivalent to
\[
q_2\le0,
\qquad
q_1+q_2\ge0~.
\]
Thus local dominance does not force $q=0$.  The core ray still rules out
negative comparison mass above the core cutoff.  Any remaining negative mass
lies above the redundant cutoff, which yields $q_2<0$.  This is consistent
with the preceding inequalities and is precisely the direction used by the
gated dominator.  In the three-dimensional model, the second unrestricted
core supplies both $(1,s,0)$ and $(s,1,0)$ for every real $s$.  These
directions collapse the permissible first-moment cone to $\{0\}$, converting
the same tail comparison into a contradiction.
\end{remark}

\section{One additional supporting halfspace}

The rectangle proof has a broader form.  The zero-mass and first-moment
arguments in that proof use only the two-core model and exact size, not the
shape of the acceptance region.  Thus, for any nice $A$ and any weak
dominator $\psi$, with $h:=\psi-\phi_A$,
\begin{equation}\label{eq:model-wide-balance}
E_0h=0,
\qquad
E_0[Xh(X)]=0~.
\end{equation}
The first identity follows by approaching the origin along $t e_1$.  For the
second, differentiation along $(1,s,0)$ and $(s,1,0)$ for every
$s\in\mathbf R$, followed by $(1,0,1)$ and $(1,0,-1)$, forces all three
components to vanish.

\begin{proposition}[Feasible support plus one halfspace]
\label{prop:one-extra-halfspace}
Let
\[
A
=
A_g\cap\{x:u'x\le c\}
\]
be nice, where $u\ne0$ is arbitrary and $A_g$ is an intersection of
halfspaces with feasible normals.  Then $\phi_A$ is $\alpha$-admissible.
\end{proposition}

\begin{proof}
Let $\psi$ weakly dominate $\phi_A$ and define $h:=\psi-\phi_A$.
For every feasible halfspace defining $A_g$, the exponential-tilt argument
in Proposition~\ref{prop:feasible-support} eliminates $h^-$ beyond that
halfspace.  A countable-subcover argument therefore gives
\[
h^-=0\quad\text{almost surely on }A_g^c~.
\]
Since $h^-$ is supported on $A^c$, any remaining negative part is supported
on $A_g\cap\{u'X>c\}$, while $h^+$ is supported on
$A\subseteq\{u'X\le c\}$.  By~\eqref{eq:model-wide-balance},
$E_0h^+=E_0h^-=:M$ and $E_0[u'Xh]=0$.  If $M>0$, then
\[
E_0[u'Xh^+]\le cM,
\qquad
E_0[u'Xh^-]>cM~,
\]
so $E_0[u'Xh]<0$, a contradiction.  Hence $M=0$, and mutual absolute
continuity finishes the proof.
\end{proof}

Proposition~\ref{prop:one-extra-halfspace} contains the rectangle result as
a special case: take
\[
A_g=\{x_1\le c_1,\ x_2\le c_2\},
\qquad
u=e_3~.
\]
It also covers a halfspace using only the redundant coordinate by taking
$A_g=\mathbf R^3$.

The qualification ``one halfspace'' is substantive.  The proposition does
not show that a rectangle with an additional bad cut is admissible: the
rectangle already uses the bad normal $e_3$, so the additional bad cut would
create two bad facets.  More generally, any inadmissible nice polyhedral test
must require at least two nonproportional bad facets and cannot be rewritten
using feasible halfspaces plus a single additional halfspace.

A minimal unresolved family is
\begin{equation}\label{eq:two-bad-facet-family}
A(\lambda)
:=
\{x_1\le\kappa,\ x_2\le\kappa,\
x_1+\lambda x_3\le b,\ x_2+\lambda x_3\le b\},
\qquad \lambda>1~,
\end{equation}
with $(\kappa,b)$ calibrated to exact size.  Its two gated-face normals
$(1,0,\lambda)$ and $(0,1,\lambda)$ are bad and nonproportional.  Neither
Proposition~\ref{prop:feasible-support} nor
Proposition~\ref{prop:one-extra-halfspace} determines its admissibility.
At $\lambda=1$ the family is the admissible symmetric gate, while after
appropriate rescaling it approaches a rectangle as $\lambda\to\infty$.
Whether some intermediate member is inadmissible remains open.  Thus the
conjecture that every nice test is admissible is plausible, but is not a
consequence of the results above.

\section{A redundant-only test and an inadmissible example}

Proposition~\ref{prop:one-extra-halfspace} immediately handles the halfspace
that uses only the redundant coordinate.  We record this special case
explicitly.  Its admissibility does not imply that every exact-size test is
admissible in the present model.

Let $z_\alpha:=\Phi^{-1}(1-\alpha)$ and define
\begin{equation}\label{eq:x3-only-test}
\phi_3(x):=I\{x_3>z_\alpha\}~.
\end{equation}
Since every null mean satisfies $\mu_3\le0$, this test has level $\alpha$,
with size attained at the origin.

\begin{proposition}[Redundant-only halfspace]
\label{prop:x3-only-admissible}
Suppose $0<\alpha<1$.  The test $\phi_3$ in~\eqref{eq:x3-only-test} is
$\alpha$-admissible against all level-$\alpha$ tests, including randomized
competitors.
\end{proposition}

\begin{proof}
Apply Proposition~\ref{prop:one-extra-halfspace} with
$A_g=\mathbf R^3$, $u=e_3$, and $c=z_\alpha$.
\end{proof}

The exact calibration in~\eqref{eq:x3-only-test} matters.  If the cutoff is
larger than $z_\alpha$, the resulting conservative test is dominated by
$\phi_3$; if it is smaller, the test is not level $\alpha$.

There are nonetheless inadmissible exact-size tests.  The following example
uses all three coordinates but allocates rejection probability to evidence
that the three means are equal.

Define
\begin{equation}\label{eq:contrast-vector}
Y
:=
\begin{pmatrix}
X_3-X_1\\
X_3-X_2
\end{pmatrix},
\qquad
V
:=
\begin{pmatrix}
2&1\\
1&2
\end{pmatrix}~.
\end{equation}
Under $P_0$, $Y\sim N_2(0,V)$.  Let $r_\alpha$ denote the lower
$\alpha$-quantile of a $\chi^2_2$ distribution and consider
\begin{equation}\label{eq:bad-contrast-test}
\phi_{\mathrm{eq}}(x)
:=
I\{Y'V^{-1}Y\le r_\alpha\}~.
\end{equation}

\begin{proposition}[An inadmissible exact-size test]
\label{prop:inadmissible-contrast-test}
The test $\phi_{\mathrm{eq}}$ is a nonrandomized level-$\alpha$ test with
size attained at the origin, but it is inadmissible against randomized
competitors.
\end{proposition}

\begin{proof}
At a general mean, let
\[
\delta
:=
\begin{pmatrix}
\mu_3-\mu_1\\
\mu_3-\mu_2
\end{pmatrix}~.
\]
Then $Y'V^{-1}Y$ has a noncentral $\chi^2_2$ distribution with noncentrality
parameter $\delta'V^{-1}\delta$.  Its probability of falling below the fixed
positive threshold $r_\alpha$ is at most its probability under zero
noncentrality.  Consequently,
\[
E_\mu\phi_{\mathrm{eq}}\le\alpha
\qquad\text{for every }\mu\in\mathbf R^3~,
\]
with equality when $\delta=0$ and strict inequality when $\delta\ne0$.
The constant randomized test $\psi(x)\equiv\alpha$ is level $\alpha$ and
has power $\alpha$ at every mean.  It therefore weakly dominates
$\phi_{\mathrm{eq}}$, with strict inequality, for example, at
$\mu=(1,0,0)\in\Lambda_1$.
\end{proof}

The test in~\eqref{eq:bad-contrast-test} is deliberately outside the natural
class considered earlier: it is not coordinatewise monotone, and its
acceptance region is not convex.  We do not currently have an inadmissible
example among exact-size tests with closed, convex, coordinatewise-lower
acceptance regions in the two-core model.  The results above establish
admissibility for nonnegative halfspaces, finite rectangles, feasibly
supported regions, and regions obtained from feasible support by one
additional halfspace.  The two-bad-facet family in
\eqref{eq:two-bad-facet-family} remains unresolved.

\section{One core and two redundant moments}

Consider now the alternative partition
\[
\mathcal C=\{1\},
\qquad
\mathcal R=\{2,3\}~.
\]
The maintained mean space becomes
\begin{equation}\label{eq:one-core-space}
\Lambda^{1\mid23}
:=
\{\mu\in\mathbf R^3:\mu_2\le\mu_1,\ \mu_3\le\mu_1\}~.
\end{equation}
Its null and alternative are
\[
\Lambda^{1\mid23}_0
:=
\Lambda^{1\mid23}\cap(-\infty,0]^3,
\qquad
\Lambda^{1\mid23}_1
:=
\{\mu\in\Lambda^{1\mid23}:\mu_1>0\}~.
\]
Unlike the two-core model, the same core coordinate is positive and dominates
both redundant means at every alternative.  This is exactly the condition
needed by the bivariate gated transfer.

Return to the finite rectangle $A_c$ in~\eqref{eq:general-rectangle}, with
$P_0(A_c)=1-\alpha$.  Fix $r\in\{2,3\}$, and let $s$ denote the other
redundant coordinate.  Choose the unique finite $b_r$ satisfying
\begin{equation}\label{eq:one-core-gate-calibration}
P_0\{X_1\le c_1,\ X_1+X_r\le b_r\}
=
\Phi(c_1)\Phi(c_r)~.
\end{equation}
Existence and uniqueness follow because the left side is continuous and
strictly increasing in $b_r$, from $0$ to $\Phi(c_1)$, while the right side
lies strictly between these limits.
Define the gated acceptance region
\begin{equation}\label{eq:one-core-gated-rectangle}
A_c^{(r)}
:=
\{x_1\le c_1,\ x_1+x_r\le b_r,\ x_s\le c_s\}~.
\end{equation}

\begin{proposition}[Rectangles are inadmissible with one core]
\label{prop:one-core-rectangle-dominated}
For either $r\in\{2,3\}$, the test with acceptance region $A_c^{(r)}$ has
exact size $\alpha$ and strictly dominates the rectangular test with
acceptance region $A_c$ over $\Lambda^{1\mid23}_1$.  Hence every finite
exact-size rectangle, including the equal-cutoff max test, is inadmissible.
\end{proposition}

\begin{proof}
Independence and~\eqref{eq:one-core-gate-calibration} give
$P_0(A_c^{(r)})=P_0(A_c)=1-\alpha$.  Coordinatewise lowerness then gives
level $\alpha$ over $\Lambda^{1\mid23}_0$.

In the $(1,r)$ coordinates, the points newly rejected by the gated test form
\[
G_r
:=
\{x_1\le c_1,\ x_r\le c_r,\ x_1+x_r>b_r\}~,
\]
whereas the points newly accepted form
\[
D_r
:=
\{x_1\le c_1,\ x_r>c_r,\ x_1+x_r\le b_r\}~.
\]
Calibration gives $P_0(G_r)=P_0(D_r)$.  This common probability is positive.
Indeed, calibration forces $b_r<c_1+c_r$: otherwise the original bivariate
rectangle would be strictly contained, up to a positive-probability set, in
the gated region, contradicting~\eqref{eq:one-core-gate-calibration}.  For
every $g\in G_r$ and
$d\in D_r$,
\[
g_1+g_r>d_1+d_r,
\qquad
-g_r>-d_r~.
\]
At any $\mu\in\Lambda^{1\mid23}_1$, write
\[
u:=\mu_1>0,
\qquad
v:=\mu_1-\mu_r\ge0~.
\]
The likelihood ratio in the selected coordinates is proportional to
\[
\exp\{u(x_1+x_r)+v(-x_r)\}~.
\]
For a direct strict comparison, set
\[
\ell(x):=u(x_1+x_r)-v x_r,
\qquad
L:=u b_r-v c_r~.
\]
Then $\ell(g)>L$ for every $g\in G_r$, whereas $\ell(d)\le L$ for every
$d\in D_r$.  Consequently,
\[
\int_{G_r}e^{\ell(x)}\,dP_0(x)
>
e^L P_0(G_r)
=
e^L P_0(D_r)
\ge
\int_{D_r}e^{\ell(x)}\,dP_0(x)~.
\]
The common likelihood-ratio normalizing factor therefore gives
$P_\mu(G_r)>P_\mu(D_r)$.  The remaining constraint $X_s\le c_s$ is
independent and common to both tests, so multiplying by
$\Phi(c_s-\mu_s)>0$ proves strict power domination.
\end{proof}

Only one redundant face needs to be gated to establish inadmissibility.  The
test in~\eqref{eq:one-core-gated-rectangle} still uses all three moments: the
other redundant coordinate remains through its original rectangular face.
The same construction works with either redundant coordinate.

\enlargethispage{4\baselineskip}

It is also easy to construct admissible tests that gate both redundant
moments.  For finite $a,b_2,b_3$, consider
\begin{equation}\label{eq:double-gated-region}
A^{\mathrm{gg}}_{a,b_2,b_3}
:=
\{x_1\le a,\ x_1+x_2\le b_2,\ x_1+x_3\le b_3\}~.
\end{equation}
The feasible-normal cone for~\eqref{eq:one-core-space} is
\[
\mathcal F^{1\mid23}
=
\{a\ge0:a_2\le a_1,\ a_3\le a_1\}~.
\]
Every finite normal of~\eqref{eq:double-gated-region} is a nonnegative
combination of
\[
e_1,
\qquad
(1,1,0),
\qquad
(1,0,1)~,
\]
and therefore belongs to $\mathcal F^{1\mid23}$.  Once the three cutoffs are
calibrated to give $P_0(A^{\mathrm{gg}}_{a,b_2,b_3})=1-\alpha$, the
feasible-normal sufficient condition establishes admissibility.

The pairwise transfer proof does not show that a jointly calibrated
double-gated test dominates the original rectangle.  After the first gate is
introduced, the remaining constraints no longer factor into an independent
bivariate problem.  A randomized mixture of the two single-gate dominators
does dominate the rectangle, but deterministic domination by a double-gated
test remains unresolved.

\end{document}
